\section{The problem and the mechanism}
An experiment does not supply a geometric parameter space together with a free observation map. It supplies a preparation, physical commands, causal reports, and executable future tests. A visible history determines predictions of those tests. The resulting predictive quotient may be infinite, singular, or model dependent. Even when the terminal task is finite dimensional, a favorable task value need not be obtainable at a prescribed acquisition cost or executable with a prescribed finite controller.

The problem considered here is to connect these objects in their physical order:
\[
 \begin{gathered}
 \text{prepared kernel}\longrightarrow\text{future tests}
 \longrightarrow\text{predictive quotient}\\
 \longrightarrow\text{causal acquisition}
 \longrightarrow\text{acquired geometry}\longrightarrow\text{risk}.
 \end{gathered}
\]
Three distinctions are essential. A law obtained by selecting favorable reports is not a free preparation. A task proxy is not necessarily the full predictive state. The labels used to acquire a report are not available a second time as terminal prediction values.

A finite observable projection can make exact stopped occupation tractable, but it is restrictive: two histories with the same protocol phase may induce different next-report laws. We replace this exact closure by two inequalities verified before an observer or a budget is chosen. The first is a predictable lower certificate for the physical effort needed to obtain productive reports. The second is a uniform upper certificate for an explicitly executable acquisition program. Both allow the microscopic state, the predictive coordinate and the report law to retain the entire relevant past. Conditional comparison with one actually acquired reference law supplies the remaining geometric control.

Theorem~\ref{thm:main} is the central result. It bounds the risks of all legal history-dependent procedures below, and constructs one parameter-independent finite-state procedure for the upper bound. Its two profiles have different, explicit budget rescalings. They become matching curves when the acquired geometry has the stated regularity; they are not silently identified at an atomic critical point. Theorem~\ref{thm:herald} applies the theorem to genuinely state-dependent heralding, with no independence of successive attempts and no common finite report row. The unknown heralding rate may change with hidden state and history inside certified bounds. Theorem~\ref{thm:gaussian-herald} also derives the ordinary hidden-sign risk without a supplied reliability coordinate, giving a stretched-logarithmic acquisition law for endogenous waiting.

A second issue concerns the origin of the task geometry. Theorem~\ref{thm:fiber} gives a necessary and sufficient criterion for a \emph{uniform quadratic mark of the declared kind}. On each actual full-information task fiber, form the least residual variance and the minimax quadratic radius. Such a mark exists precisely when the radius is bounded by a fixed multiple of that least variance. This is not a criterion for every possible compression method, but it identifies exactly the obstruction to the conditional comparison used here. It also makes equivalent marks interchangeable up to fixed constants.

The physical controller is handled in Section~\ref{sec:memory}. Commands are identified by their physical interface, including legality, not merely by their printed names. Mandatory incompatible physical command sets give a live-state lower bound. This bound, and the optimal live width itself, are invariant under stateless resource-preserving recodings. A recoding that needs to remember a protocol phase is not stateless and its state is charged. Appendix~\ref{sec:flow} separately extends terminal-submeasure-preserving killed-flow realization from finite drivers to standard Borel state spaces. That realization may require the full predictive state; it is never presented as a free finite-memory implementation.

Conditional-envelope rejection was already used in the causal-stopping and causal-acquisition companions~\cite{GTF46,GTF51}; it is not claimed as a new sampling method here. The marked quantization calculation and conditional projection mechanism are inherited from the stopped-experiment and calibrated-protocol companions~\cite{GTF52,GTF53} and are re-proved where used. Classical occupation-measure Markovization is adjacent to, and not invented by, the present work: finite extreme occupations in Borel absorbing models are studied by Piunovskiy and Zhang~\cite{PZExtreme}; strategic-to-occupation continuity and absorption by the same authors~\cite{PZContinuity}; and stationary or chattering sufficiency by Dufour and Prieto-Rumeau~\cite{DPAbsorb,DPChatter}. Our main risk proof does not rely on compactness of the space of all stopped trajectories. It uses predictable count domination and an explicit finite-state acquisition construction.

Finite observer memory alone does not make general partially observed stopping easy. Optimal-cost almost-sure reachability~\cite{CCGK} and fixed-memory reachability~\cite{ACSS} have quite different computational difficulties. The theorem below uses substantive task, mass and effort certificates; it is not an efficient solution of arbitrary POMDPs. The new point relative to the companions is the combination of raw-history-dependent acquisition, a sharp calibration criterion, physical-command accounting, and fully paid implementation. No claim of an exhaustive priority classification or of a solution to every causal experiment is made.
